\documentclass[../main.tex]{subfiles}

\begin{document}

\section*{第1問}

\begin{proof}[解]
反射光が半円の中心$O$を中心とする単位円周上を進む様子を考える。$k$回目の反射点を$Q_k$とする。円の反射では入射角と反射角（半径となす角）が等しいことから，各反射点の間で弧の中心角は常に一定値$(\pi-2\theta)$ずつ進む（$A=(1,0)$から出発し，$Q_1$の中心角は$\pi-2\theta$）。したがって
\[
Q_1:\ \pi-2\theta,\qquad Q_2:\ 2(\pi-2\theta)=2\pi-4\theta,\qquad Q_3:\ 3(\pi-2\theta)
\]
（角は正の$x$軸から反時計回りに測った偏角）。

\textbf{(1)} 2回目の反射（$Q_2$）が起こり，かつ3回目の反射（$Q_3$）が起こらない条件を求める。$Q_1,Q_2$が上半円周上にあるためには$0<\theta<\pi/2\ \cdots\text{①}$が必要（$Q_1$の偏角$\pi-2\theta\in(0,\pi)$から）。$Q_2$が有効な反射点であるためには
\[
0<2\pi-4\theta<\pi
\]
かつ，3回目の反射が起こらない（$Q_3$の偏角が$\pi$を超え，光線がその前に$x$軸を通過する）ためには
\[
\pi<3(\pi-2\theta)
\]
これらをまとめて
\[
0<2(\pi-2\theta)<\pi<3(\pi-2\theta) \quad(\because\text{①})
\]
第1不等式$2(\pi-2\theta)<\pi$から$\theta>\pi/4$，第2不等式$\pi<3(\pi-2\theta)$から$\theta<\pi/3$。よって
\[
\frac\pi4<\theta<\frac\pi3
\]

\textbf{(2)} $Q_2$での反射角も$\theta$（半径$OQ_2$との角）であり，反射光$Q_2P$は$x$軸上の点$P$へ向かう。$OQ_2$の偏角は$2\pi-4\theta\in(2\pi/3,\pi)$（①，(1)の範囲で）で，$OP$は負の$x$軸方向（偏角$\pi$）だから，$\triangle OQ_2P$において
\[
\angle OQ_2P=\theta,\qquad \angle Q_2OP=\pi-(2\pi-4\theta)=4\theta-\pi,\qquad OQ_2=1
\]
角の和から$\angle OPQ_2=\pi-\theta-(4\theta-\pi)=2\pi-5\theta$。正弦定理$\dfrac{OP}{\sin\theta}=\dfrac{OQ_2}{\sin\angle OPQ_2}$より
\[
OP=\frac{\sin\theta}{\sin(2\pi-5\theta)}=\frac{\sin\theta}{-\sin5\theta}
\]
$P$は負の$x$軸上（$P_x=-OP$）だから
\[
P=\left(\frac{\sin\theta}{\sin5\theta},\ 0\right)
\]

\textbf{(3)} $f(\theta)=\dfrac{\sin\theta}{\sin5\theta}$とおく。以下$\sin\theta=s,\ \cos\theta=c$と書く。
\[
\sin5\theta=\sin(2\theta+3\theta)=\sin2\theta\cos3\theta+\sin3\theta\cos2\theta
\]
\[
=2sc(4c^3-3c)+(3s-4s^3)(1-2s^2)=2s(1-s^2)(1-4s^2)+(3s-4s^3)(1-2s^2)
\]
\[
=(2s-10s^3+8s^5)+(3s-10s^3+8s^5)=5s-20s^3+16s^5=s(16s^4-20s^2+5)
\]
$\dfrac\pi4<\theta<\dfrac\pi3$では$s\ne0$だから
\[
f(\theta)=\frac1{16s^4-20s^2+5}
\]
$t=s^2$とおくと，$\theta\in(\pi/4,\pi/3)$で$\sin\theta$は単調増加だから$t\in(1/2,3/4)$。
\[
f=\frac1{16t^2-20t+5}=\frac1{16\left(t-\frac58\right)^2-\frac54}
\]
$g(t)=16(t-5/8)^2-5/4$は$t=5/8\in(1/2,3/4)$で最小値$-5/4$をとり，$t\to1/2^+,\ 3/4^-$でともに$g\to-1$（$g(1/2)=g(3/4)=-1$，開区間なので到達しない）。よって$g(t)\in[-5/4,-1)$（負値）。したがって$f=1/g$は$g=-5/4$で$f=-4/5$（達成），$g\to-1^-$で$f\to-1$（到達しない）となり
\[
-1<f(\theta)\le-\frac45
\]
以上から，$P$の動く範囲は，$x$軸上の
\[
-1<x\le-\frac45
\]
（$y=0$）である。
\end{proof}

\newpage

\section*{第2問}

\begin{proof}[解]
$\cos\theta=c,\ \sin\theta=s$と書く。

\textbf{(1)}
\[
\left|z+\frac12\right|<\frac12 \iff \left|(rc+\tfrac12)+irs\right|<\frac12 \iff \left(rc+\frac12\right)^2+(rs)^2<\frac14 \quad(\because\text{両辺}\ge0)
\]
\[
\iff r^2+rc<0 \iff r(r+\cos\theta)<0
\]
$r\ge0$だから，これは$r\ne0$かつ$r+\cos\theta<0$と同値。すなわち
\[
r>0\ \text{かつ}\ r<-\cos\theta \quad(\text{したがって}\ \cos\theta<0\ \text{も必要})
\]

\textbf{(2)} $z\ne1$の時，$1+z+\cdots+z^n=\dfrac{z^{n+1}-1}{z-1}$だから
\[
\left|1+z+\cdots+z^n\right|^2=\frac{|z^{n+1}-1|^2}{|z-1|^2}
\]
$|z-1|^2=(rc-1)^2+(rs)^2=r^2-2rc+1=r^2-2r\cos\theta+1$。
\[
|z^{n+1}-1|^2=\left(r^{n+1}\cos(n+1)\theta-1\right)^2+\left(r^{n+1}\sin(n+1)\theta\right)^2=r^{2n+2}-2r^{n+1}\cos(n+1)\theta+1
\]
よって
\[
|1+z+\cdots+z^n|^2=\frac{r^{2n+2}-2r^{n+1}\cos(n+1)\theta+1}{r^2-2r\cos\theta+1}
\]
（$z=1$（$r=1,\theta=0$）の時は直接$|1+z+\cdots+z^n|^2=(n+1)^2$。）

\textbf{(3)} $Z_n=1+z+\cdots+z^n$とおく。(1)の条件から$0<r<-\cos\theta$，特に$\cos\theta<0$。まず
\[
r^2-2r\cos\theta+1=(r-1)^2+2r(1-\cos\theta)>0 \quad(\because\cos\theta<0\ \text{より}\ 1-\cos\theta>0,\ r>0)
\]
だから(2)の分母は正。$|Z_n|<1\iff|Z_n|^2<1\iff$（分母を払って）
\[
r^{2n+2}-2r^{n+1}\cos(n+1)\theta+1<r^2-2r\cos\theta+1
\]
\[
\iff r^{2n+2}-2r^{n+1}\cos(n+1)\theta<r^2-2r\cos\theta
\]
両辺を$r>0$で割って，示すべき式は
\[
r^{2n+1}-2r^n\cos(n+1)\theta<r-2\cos\theta \quad\cdots\text{①}
\]
これを示す。$r+\cos\theta<0$（(1)）から$2\cos\theta<-2r$，よって
\[
r^{2n+1}-2r^n\cos(n+1)\theta-r+2\cos\theta<r^{2n+1}-2r^n\cos(n+1)\theta-r-2r
\]
\[
=r\left(r^{2n}-2r^{n-1}\cos(n+1)\theta-3\right)
\]
$|\cos(n+1)\theta|\le1$より$-\cos(n+1)\theta\le1$だから
\[
<r\left(r^{2n}+2r^{n-1}-3\right)
\]
さらに$0<r<-\cos\theta\le1$から$0<r<1$，よって$r^{2n}<1,\ r^{n-1}\le1$（$n\ge1$）だから
\[
r^{2n}+2r^{n-1}-3<1+2-3=0
\]
以上から
\[
r^{2n+1}-2r^n\cos(n+1)\theta-r+2\cos\theta<r\cdot(\text{負})<0
\]
すなわち①が成り立ち，したがって$|Z_n|^2<1$，すなわち$|1+z+\cdots+z^n|<1$が全ての自然数$n$で成り立つ。
\end{proof}

\newpage

\section*{第3問}

\begin{proof}[解]
三角柱の底面（正三角形，1辺1）を$DEF$，対応する上面を$ABC$（$AD=BE=CF=2$，各々垂直な辺）とする。断面の直角三角形$PQR$の頂点は辺$AD,BE,CF$上に1つずつあり，対称性からそのうちの1つが$E$（底面の頂点）に一致する場合を考えれば良い。$P=E$，$Q$が辺$AD$上，$R$が辺$CF$上にあるとし，
\[
DQ=x,\qquad FR=x+t \qquad(0\le x\le x+t\le2)\ \cdots(\ast)
\]
とおく（$D,E,F$は底面の3頂点で互いの距離は1）。

$P=E$は高さ$0$，$Q$は高さ$x$，$R$は高さ$x+t$にあり，$D,E,F$間の水平距離はいずれも$1$だから
\[
PQ=\sqrt{1+x^2},\qquad QR=\sqrt{1+t^2},\qquad PR=\sqrt{1+(x+t)^2}
\]
$(\ast)$より$PR$が最大辺だから，$PR$が斜辺，すなわち直角は$Q$にある。$\triangle PQR$の面積$S$は
\[
S=\frac12\,QP\cdot QR=\frac12\sqrt{(1+x^2)(1+t^2)} \quad\cdots\text{①}
\]
また，ピタゴラスの定理$PQ^2+QR^2=PR^2$から
\[
(1+x^2)+(1+t^2)=1+(x+t)^2 \iff 2+x^2+t^2=1+x^2+2xt+t^2 \iff 2xt=1
\]
\[
\therefore\ xt=\frac12 \quad\cdots\text{②}
\]

$(\ast)$，②のもとで①の値域をもとめる。$k=x+t$とおくと，②から$x,t$は2次方程式
\[
p^2-kp+\frac12=0 \quad(0\le k\le2)
\]
の2つの正の実解であり（$xt=1/2>0$から$x,t$は同符号，かつ$k\ge0$から正），判別式$D=k^2-2\ge0$が必要十分。$0\le k\le2$とあわせて
\[
\sqrt2\le k\le2 \quad\cdots\text{③}
\]
①より
\[
S=\frac12\sqrt{1+x^2+t^2+x^2t^2}=\frac12\sqrt{1+\{(x+t)^2-2xt\}+x^2t^2}=\frac12\sqrt{1+(k^2-1)+\frac14}=\frac12\sqrt{k^2+\frac14}
\]
これは$k>0$で単調増加だから，③より$k=\sqrt2$で最小，$k=2$で最大：
\[
S_{\min}=\frac12\sqrt{2+\frac14}=\frac12\cdot\frac32=\frac34,\qquad
S_{\max}=\frac12\sqrt{4+\frac14}=\frac12\cdot\frac{\sqrt{17}}2=\frac{\sqrt{17}}4
\]
したがって，もとめる面積の範囲は
\[
\frac34\le S\le\frac{\sqrt{17}}4
\]
\end{proof}

\newpage

\section*{第4問}

\begin{proof}[解]
$f(x)=e^{nx}-1-e^x$，$t=e^x$とおく。

\textbf{(1)} $f'(x)=ne^{nx}-e^x=nt^n-t=t(nt^{n-1}-1)$。$x\ge0$では$t=e^x\ge1$だから，$n\ge2$より$nt^{n-1}\ge n\ge2>1$，よって$nt^{n-1}-1>0$，$t>0$とあわせて$f'(x)>0$（$x\ge0$）。
\begin{center}
\begin{tabular}{c|cc}
$x$ & $0$ & \\
\hline
$f'$ & & $+$ \\
\hline
$f$ & & $\nearrow$
\end{tabular}
\end{center}
これと$f(0)=1-1-1=-1<0$，$f(x)\to\infty\ (x\to\infty)$から，$f(x)=0$は$x>0$にただ1つの実解を持つ（$f$は$x\ge0$で狭義単調増加だから一意）。したがって(ア)，(イ)は第1象限にただ1つの交点を持つ。

\textbf{(2)} 交点の$x$座標を$a_n$とする（$f(a_n)=0,\ a_n>0$）。
\[
f\left(\frac1n\right)=e-1-e^{1/n}
\]
$n\ge2$より$1/n\le1/2$，よって$e^{1/n}\le e^{1/2}$，したがって
\[
f\left(\frac1n\right)=e-1-e^{1/n}\ge e-1-e^{1/2}
\]
$(1.7)^2=2.89>e$より$e^{1/2}<1.7$，また$e>2.7$だから
\[
e-1-e^{1/2}>2.7-1-1.7=0
\]
よって$f(1/n)\ge e-1-e^{1/2}>0$。(1)の単調性と$f(a_n)=0$から$a_n<1/n$，$a_n>0$とあわせて
\[
0<a_n<\frac1n
\]
はさみうちの定理から$a_n\to0\ (n\to\infty)$。

また，$f(a_n)=0$より$e^{na_n}=e^{a_n}+1$。両辺の自然対数をとって
\[
na_n=\log(e^{a_n}+1) \xrightarrow{n\to\infty}\log2 \quad(\because a_n\to0)
\]

\textbf{(3)} 第1象限で(ア)が(イ)より上（$0<x<a_n$で$e^x>e^{nx}-1$）だから，
\[
S_n=\int_0^{a_n}\left\{e^x-(e^{nx}-1)\right\}dx=\left[e^x-\frac1ne^{nx}+x\right]_0^{a_n}
=e^{a_n}-\frac1ne^{na_n}+a_n-1+\frac1n
\]
\[
\therefore\ nS_n=n(e^{a_n}-1)+na_n+1-e^{na_n}=na_n\cdot\frac{e^{a_n}-1}{a_n}+na_n+1-e^{na_n}
\]
$n\to\infty$の時，$a_n\to0$より$\dfrac{e^{a_n}-1}{a_n}\to1$，$na_n\to\log2$，$e^{na_n}\to e^{\log2}=2$だから
\[
nS_n\longrightarrow(\log2)\cdot1+\log2+1-2=2\log2-1
\]
\end{proof}

\end{document}
